% GENERATED FILE. Edit canonical lessons, then run npm run generate:books.
% canonical-source-hash: fnv1a64:6913359837475824
% canonical-lessons: TA-C493-anni, TA-C493-koluntan, TA-C493-nattanar, TA-C493-campanti, TA-C493-citti

\chapter{In-Laws}
\label{ch:ta-in-laws}
\hlchaptermodality{493}

\begin{chapteropening}
\textbf{By the end of this chapter:} \emph{I can name in-laws on both sides of a marriage.}

Complete the last lesson of chapter 493: I can name in-laws on both sides of a marriage.
\end{chapteropening}

\section[\texorpdfstring{\ta{அண்ணி} (aṇṇi)}{அண்ணி (aṇṇi)}]{\ta{அண்ணி} (aṇṇi) — an elder brother’s wife}

\label{lesson:TA-C493-anni}



\begin{quote}
Before the new one: say the Tamil for a sister, then the Tamil for the youngest child.
\end{quote}

\subsection*{You'll want to know: \ta{அண்ணி}}
\textbf{\ta{அண்ணி}} — \emph{aṇṇi} — \textquotedblleft{}an elder brother’s wife\textquotedblright{}.

\begin{quote}
\ta{என்} \ta{அண்ணி} \ta{நன்றாகச்} \ta{சமைப்பாள்}.
\end{quote}

\begin{quote}
\emph{eṉ aṇṇi naṉṟākac camaippāḷ.}
\end{quote}

\begin{quote}
My elder brother’s wife cooks well.
\end{quote}

Made from \ta{அண்ணன்} \textquotedblleft{}elder brother\textquotedblright{} with the -\ta{இ} of a woman.

\begin{cousinweb}
\textbf{\ta{அண்ணி}} (\emph{aṇṇi}) is native Dravidian, built on \ta{அண்ணன்} \emph{aṇṇaṉ}, \textquotedblleft{}elder brother\textquotedblright{} (DEDR 131).
\end{cousinweb}

\begin{grammarlens}[title={What you've built}]
One more word for describing what you lived and felt.
\end{grammarlens}

\subsection*{Your turn}
\begin{itemize}
\raggedright
  \item \emph{Say it:} \emph{aṇṇi}
  \item \emph{Say it:} \emph{aṇṇi}, once more
  \item \emph{Say it:} \emph{eṉ aṇṇi naṉṟākac camaippāḷ.} — My elder brother’s wife cooks well.
  \item \emph{Recall:} say the Tamil for a sister, then the Tamil for the youngest child, then say \emph{aṇṇi} again
\end{itemize}

\begin{tcolorbox}[breakable,colback=teal!4,colframe=teal!35!black,title={Before you move on}]
What does \ta{அண்ணி} mean? (\textquotedblleft{}An elder brother’s wife\textquotedblright{}.) Say it once more.
\end{tcolorbox}
\section[\texorpdfstring{\ta{கொழுந்தன்} (koḻuntaṉ)}{கொழுந்தன் (koḻuntaṉ)}]{\ta{கொழுந்தன்} (koḻuntaṉ) — a husband’s younger brother}

\label{lesson:TA-C493-koluntan}



\begin{quote}
Before the new one: say the Tamil for the youngest child, then the Tamil for an elder brother’s wife.
\end{quote}

\subsection*{You'll want to know: \ta{கொழுந்தன்}}
\textbf{\ta{கொழுந்தன்}} — \emph{koḻuntaṉ} — \textquotedblleft{}a husband’s younger brother\textquotedblright{}.

\begin{quote}
\ta{என்} \ta{அக்காவின்} \ta{கொழுந்தன்} \ta{ஒரு} \ta{மருத்துவர்}.
\end{quote}

\begin{quote}
\emph{eṉ akkāviṉ koḻuntaṉ oru maruttuvar.}
\end{quote}

\begin{quote}
My elder sister’s brother-in-law is a doctor.
\end{quote}

From a woman’s side: her husband’s younger brother.

\begin{cousinweb}
\textbf{\ta{கொழுந்தன்}} (\emph{koḻuntaṉ}) is native Dravidian (DEDR 2149), in the entry of \ta{கொழுந்து} \emph{koḻuntu}, \textquotedblleft{}a tender shoot\textquotedblright{}.
\end{cousinweb}

\begin{grammarlens}[title={What you've built}]
One more word for describing what you lived and felt.
\end{grammarlens}

\subsection*{Your turn}
\begin{itemize}
\raggedright
  \item \emph{Say it:} \emph{koḻuntaṉ}
  \item \emph{Say it:} \emph{koḻuntaṉ}, once more
  \item \emph{Say it:} \emph{eṉ akkāviṉ koḻuntaṉ oru maruttuvar.} — My elder sister’s brother-in-law is a doctor.
  \item \emph{Recall:} say the Tamil for the youngest child, then the Tamil for an elder brother’s wife, then say \emph{koḻuntaṉ} again
\end{itemize}

\begin{tcolorbox}[breakable,colback=teal!4,colframe=teal!35!black,title={Before you move on}]
What does \ta{கொழுந்தன்} mean? (\textquotedblleft{}A husband’s younger brother\textquotedblright{}.) Say it once more.
\end{tcolorbox}
\section[\texorpdfstring{\ta{நாத்தனார்} (nāttaṉār)}{நாத்தனார் (nāttaṉār)}]{\ta{நாத்தனார்} (nāttaṉār) — a husband’s sister}

\label{lesson:TA-C493-nattanar}



\begin{quote}
Before the new one: say the Tamil for an elder brother’s wife, then the Tamil for a husband’s younger brother.
\end{quote}

\subsection*{You'll want to know: \ta{நாத்தனார்}}
\textbf{\ta{நாத்தனார்}} — \emph{nāttaṉār} — \textquotedblleft{}a husband’s sister\textquotedblright{}.

\begin{quote}
\ta{அவளும்} \ta{அவளுடைய} \ta{நாத்தனாரும்} \ta{நல்ல} \ta{நண்பர்கள்}.
\end{quote}

\begin{quote}
\emph{avaḷum avaḷuṭaiya nāttaṉārum nalla naṇparkaḷ.}
\end{quote}

\begin{quote}
She and her husband’s sister are good friends.
\end{quote}

Respectful -\ta{ஆர்} ending.

\begin{cousinweb}
\textbf{\ta{நாத்தனார்}} (\emph{nāttaṉār}) is native Dravidian (DEDR 3644); Malayalam \emph{nāttūn} is the same word.
\end{cousinweb}

\begin{grammarlens}[title={What you've built}]
One more word for describing what you lived and felt.
\end{grammarlens}

\subsection*{Your turn}
\begin{itemize}
\raggedright
  \item \emph{Say it:} \emph{nāttaṉār}
  \item \emph{Say it:} \emph{nāttaṉār}, once more
  \item \emph{Say it:} \emph{avaḷum avaḷuṭaiya nāttaṉārum nalla naṇparkaḷ.} — She and her husband’s sister are good friends.
  \item \emph{Recall:} say the Tamil for an elder brother’s wife, then the Tamil for a husband’s younger brother, then say \emph{nāttaṉār} again
\end{itemize}

\begin{tcolorbox}[breakable,colback=teal!4,colframe=teal!35!black,title={Before you move on}]
What does \ta{நாத்தனார்} mean? (\textquotedblleft{}A husband’s sister\textquotedblright{}.) Say it once more.
\end{tcolorbox}
\section[\texorpdfstring{\ta{சம்பந்தி} (campanti)}{சம்பந்தி (campanti)}]{\ta{சம்பந்தி} (campanti) — a co-parent-in-law}

\label{lesson:TA-C493-campanti}



\begin{quote}
Before the new one: say the Tamil for a husband’s younger brother, then the Tamil for a husband’s sister.
\end{quote}

\subsection*{You'll want to know: \ta{சம்பந்தி}}
\textbf{\ta{சம்பந்தி}} — \emph{campanti} — \textquotedblleft{}a co-parent-in-law\textquotedblright{}.

\begin{quote}
\ta{திருமணத்தில்} \ta{இரண்டு} \ta{சம்பந்திகளும்} \ta{சந்தோஷமாக} \ta{இருந்தார்கள்}.
\end{quote}

\begin{quote}
\emph{tirumaṇattil iraṇṭu campantikaḷum cantōṣamāka iruntārkaḷ.}
\end{quote}

\begin{quote}
At the wedding both sets of in-laws were happy.
\end{quote}

Plural: \ta{சம்பந்திகள்}, \emph{campantikaḷ}.

Your son’s or daughter’s parents-in-law are your \ta{சம்பந்தி}.

\begin{cousinweb}
\textbf{\ta{சம்பந்தி}} (\emph{campanti}) is borrowed from Sanskrit \emph{sambandhin}, \textquotedblleft{}connected, related\textquotedblright{} (Monier-Williams).
\end{cousinweb}

\begin{grammarlens}[title={What you've built}]
One more word for describing what you lived and felt.
\end{grammarlens}

\subsection*{Your turn}
\begin{itemize}
\raggedright
  \item \emph{Say it:} \emph{campanti}
  \item \emph{Say it:} \emph{campanti}, once more
  \item \emph{Say it:} \emph{tirumaṇattil iraṇṭu campantikaḷum cantōṣamāka iruntārkaḷ.} — At the wedding both sets of in-laws were happy.
  \item \emph{Recall:} say the Tamil for a husband’s younger brother, then the Tamil for a husband’s sister, then say \emph{campanti} again
\end{itemize}

\begin{tcolorbox}[breakable,colback=teal!4,colframe=teal!35!black,title={Before you move on}]
What does \ta{சம்பந்தி} mean? (\textquotedblleft{}A co-parent-in-law\textquotedblright{}.) Say it once more.
\end{tcolorbox}
\section[\texorpdfstring{\ta{சித்தி} (citti)}{சித்தி (citti)}]{\ta{சித்தி} (citti) — a mother’s younger sister}

\label{lesson:TA-C493-citti}



\begin{quote}
Before the new one: say the Tamil for a husband’s sister, then the Tamil for a co-parent-in-law.
\end{quote}

\subsection*{You'll want to know: \ta{சித்தி}}
\textbf{\ta{சித்தி}} — \emph{citti} — \textquotedblleft{}a mother’s younger sister\textquotedblright{}.

\begin{quote}
\ta{என்} \ta{சித்தி} \ta{எங்கள்} \ta{வீட்டுக்கு} \ta{அருகில்} \ta{இருக்கிறாள்}.
\end{quote}

\begin{quote}
\emph{eṉ citti eṅkaḷ vīṭṭukku arukil irukkiṟāḷ.}
\end{quote}

\begin{quote}
My mother’s younger sister lives near our house.
\end{quote}

Also the wife of \ta{சித்தப்பா}, and a stepmother.

\begin{cousinweb}
The origin of \textbf{\ta{சித்தி}} (\emph{citti}) is not settled; it pairs with \ta{சித்தப்பா}, but our sources do not give its history.
\end{cousinweb}

\begin{grammarlens}[title={What you've built}]
One more word for describing what you lived and felt.
\end{grammarlens}

\subsection*{Your turn}
\begin{itemize}
\raggedright
  \item \emph{Say it:} \emph{citti}
  \item \emph{Say it:} \emph{citti}, once more
  \item \emph{Say it:} \emph{eṉ citti eṅkaḷ vīṭṭukku arukil irukkiṟāḷ.} — My mother’s younger sister lives near our house.
  \item \emph{Recall:} say the Tamil for a husband’s sister, then the Tamil for a co-parent-in-law, then say \emph{citti} again
\end{itemize}

\begin{tcolorbox}[breakable,colback=teal!4,colframe=teal!35!black,title={Before you move on}]
What does \ta{சித்தி} mean? (\textquotedblleft{}A mother’s younger sister\textquotedblright{}.) Say it once more.
\end{tcolorbox}
